% =====================================================================
%  Solution template -- Nestruev, "Smooth Manifolds and Observables"
%  Differential Geometry Reading Group
%
%  1. Save as solutions/chNN/<exercise>_<your-name>.tex
%     e.g. solutions/ch03/3-28_aleksandras.tex
%  2. Fill in the part marked YOUR SOLUTION at the bottom.
%  The preamble is shared by everyone: please don't change it here.
%  If you need new notation, add it under "Your own macros".
% =====================================================================
\documentclass[11pt,a4paper]{article}

% ------------------------------------------------ shared preamble (don't edit)
\usepackage[margin=2.5cm]{geometry}
\usepackage{amsmath,amssymb,amsthm,mathtools}
\usepackage{enumitem}
\usepackage{tikz-cd}
\usepackage{xcolor}
\usepackage[colorlinks=true,linkcolor=blue!50!black,urlcolor=blue!50!black]{hyperref}
\usepackage[capitalise]{cleveref}

% ---------------------------------------------------------------- metadata
\makeatletter
\newcommand{\dg@chapnum}{}
\newcommand{\dg@chaptitle}{}
\newcommand{\chapter}[2]{\renewcommand{\dg@chapnum}{#1}\renewcommand{\dg@chaptitle}{#2}}
\renewcommand{\maketitle}{%
  \begin{center}
    {\small Nestruev, \emph{Smooth Manifolds and Observables} --- reading group}\par
    \vspace{0.6em}
    {\Large\bfseries Chapter \dg@chapnum: \dg@chaptitle}\par
    \vspace{0.4em}
    {\normalsize \@author}\par
    {\small \@date}\par
  \end{center}
  \vspace{1em}}
\makeatother

% ---------------------------------------------------------------- environments
\theoremstyle{plain}
\newtheorem{theorem}{Theorem}
\newtheorem{lemma}[theorem]{Lemma}
\newtheorem{proposition}[theorem]{Proposition}
\newtheorem{corollary}[theorem]{Corollary}
\newtheorem{claim}[theorem]{Claim}
\theoremstyle{definition}
\newtheorem{definition}[theorem]{Definition}
\newtheorem{example}[theorem]{Example}
\theoremstyle{remark}
\newtheorem*{remark}{Remark}

% \begin{exercise}{3.28} ... \end{exercise}   -- restate the exercise from the book
\newenvironment{exercise}[1]
  {\par\medskip\noindent\textbf{Exercise #1.}\ \itshape}
  {\par\medskip}
% \begin{solution} ... \end{solution}
\newenvironment{solution}
  {\begin{proof}[Solution]}
  {\end{proof}}

% ---------------------------------------------------------------- notation
% Notation roughly follows Nestruev. Add to this list rather than redefining
% things in individual solution files.
\newcommand{\R}{\mathbb{R}}
\newcommand{\C}{\mathbb{C}}
\newcommand{\Z}{\mathbb{Z}}
\newcommand{\N}{\mathbb{N}}
\newcommand{\Q}{\mathbb{Q}}
\newcommand{\RP}{\mathbb{R}P}
\newcommand{\F}{\mathcal{F}}            % algebra of observables
\newcommand{\Cinf}{C^{\infty}}
\newcommand{\abs}[1]{\left\lvert #1 \right\rvert}
\newcommand{\norm}[1]{\left\lVert #1 \right\rVert}
\newcommand{\Spec}{\operatorname{Spec}}
\newcommand{\Specm}{\operatorname{Specm}}
\newcommand{\Hom}{\operatorname{Hom}}
\newcommand{\End}{\operatorname{End}}
\newcommand{\Der}{\operatorname{D}}     % derivations D(P)
\newcommand{\Diff}{\operatorname{Diff}} % differential operators Diff_k(P,Q)
\newcommand{\smbl}{\operatorname{smbl}}
\newcommand{\supp}{\operatorname{supp}}
\newcommand{\id}{\mathrm{id}}
\newcommand{\im}{\operatorname{im}}
\newcommand{\coker}{\operatorname{coker}}
\newcommand{\Ann}{\operatorname{Ann}}
\newcommand{\dd}{\,\mathrm{d}}
\newcommand{\pd}[2]{\frac{\partial #1}{\partial #2}}
% |F| : the R-spectrum (dual space) of an R-algebra F
\newcommand{\dual}[1]{\lvert #1 \rvert}

% ------------------------------------------------ your own macros (optional)


% =====================================================================
%  YOUR SOLUTION
% =====================================================================
\chapter{3}{Algebras and Points}
\author{Aleksandras Melnik}
\date{\today}

\begin{document}
\maketitle

\begin{exercise}{3.3}
Suppose that the $\R$-algebras $\F_1$ and $\F_2$, as linear spaces, are
isomorphic to the plane $\R^2 = \{(x,y)\}$. Let the multiplication in
$\F_1$ and $\F_2$ be respectively given by
\begin{align*}
  (x_1,y_1)\cdot(x_2,y_2) &= (x_1x_2,\; y_1y_2),\\
  (x_1,y_1)\cdot(x_2,y_2) &= (x_1x_2 + y_1y_2,\; x_1y_2 + x_2y_1).
\end{align*}
Find the set (manifold) $M_i$ for which the algebra $\F_i$, $i=1,2$, is
the algebra of smooth functions, explicitly indicating what function on
$M_i$ corresponds to the element $(x,y)\in\F_i$. Are the algebras $\F_1$
and $\F_2$ isomorphic?
\end{exercise}

\begin{solution}
We find the manifolds for both algebras, we could just guess them, but keeping with the spirit of the book let us find them from the $\R$-algebra homomorphisms of $\F_i$ to $\R$. But before we do that it will be helpful to set out some theory.

\begin{definition}[Basis of (a real) Vector Space ]
  For a real vector space $V$ its subset $B\subseteq V$ is called a \textit{basis} if for any $v\in V$ there exists a unique finite set of elements $b_1, \dots, b_n$ from $B$ and a unique set of real scalars $\lambda_1, \dots \lambda_n \in \R$ such that:
$$v=\lambda_1 b_1+\cdots+\lambda_n b_n $$

In other words, each vector must be expressible as a unique linear combination of the basis elements.
\end{definition}

A standard result in linear algebra is that all basis are of the same cardinality, which is called \textit{the dimension} of $V$.

\begin{definition}[Linear mapping between (real) vector spaces]
  For two real vector spaces $V_1$ and $V_2$ a \textit{linear map} (or a vector space homomorphism) is a function $f:V_1\to V_2$ such that for all $u,v\in V_1$ and $\lambda\in \R$:
\begin{itemize}
  \item $f(\lambda v)=\lambda f(v)$
  \item $f(u+v)=f(u)+f(v)$
\end{itemize}

\end{definition}

From these two definitions it follows, that a linear map $f:V_1\to V_2$ is completely determined by the values it takes on a basis $B$ of $V_1$, because for any element $v\in V_1$:

$$f(v)=f(\lambda_1 b_1+ \cdots\lambda_n b_n)=\lambda_1 f(b_1)+\cdots +\lambda_n f(b_n) $$


Now returning to the problem at hand, we are interested in $\R$-algebra homomorphisms from $\F_i$ to $\R$, which will ultimately be our ``points''. Since $\R$-algebras are linear spaces (with additional multiplicative) structure, their homomorphisms must also be linear maps. Both of the algebras at hand have the basis $B=\{(1,0), (0,1)\}$, so any homomorphism $\alpha :\F_i\to \R$ will be determined by the two values $\alpha(1,0)$ and $\alpha (0,1)$.

Furthermore, as a (unital) algebra homomorphism $\alpha$ has to preserve the multiplicative structure and the multiplicative unit, i.e.:

\begin{itemize}
  \item $\alpha(u\cdot v)= \alpha(u)\cdot \alpha(v)$
  \item $\alpha(\mathbf{1})=1$
\end{itemize}

Having made this setup we can go into the specific examples.
\begin{enumerate}
  \item First of all the multiplicative unit of $\F_1$ is evidently $(1,1)$, since $(1,1)\cdot (x,y)=(x,y)$ which means that $\alpha(1,1)=1$, and hence, by linearity, $\alpha(1,0)+\alpha(0,1)=1$. Furthermore since $\alpha$ has to preserve multiplication we have that 
    $$\alpha(1,0)\cdot \alpha(0,1)=\alpha((1,0)\cdot (0,1))=\alpha(0,0)=0$$
  Where the last equality follows from $\alpha$ being a linear map. These two facts above give the following constraints on $\alpha$:

  \begin{itemize}
    \item $\alpha(1,0)+\alpha(0,1)=1$
    \item $\alpha(1,0)\cdot\alpha(0,1)=0$
  \end{itemize}

  Since these values are real numbers, the only solutions are those where one basis vector takes the value 0 and the other takes the value 1. Indeed, these are the maps:
    \begin{itemize}
      \item $\alpha_1(x,y):=x$. This is a homomorphism since it is a linear map, i.e. it preserves the vectorial structure, and for the multiplicative structure we have:
        $$
          \alpha_1((x_1,y_1)\cdot(x_2,y_2)) = \alpha_1(x_1 x_2,y_1 y_2)=x_1x_2= \alpha_1(x_1,y_1)
          \cdot\alpha_1(x_2,y_2)$$ 
        
      It also preserves the multiplicative identity $\alpha_1(1,1)=1$
    \item The same goes for $\alpha_2(x,y):=y$.


    \end{itemize}

    Thus we have identified the two ``points'' $\alpha_1$ and $\alpha_2$ for the ``manifold'' $M_1=\{\alpha_1, \alpha_2\}$, and the function corresponding to $(x,y)\in \F_1$ is defined by:

    \begin{itemize}
      \item $\alpha_1 \overset{(x,y)}{\mapsto} x$
      \item $\alpha_2  \overset{(x,y)}{\mapsto} y$
    \end{itemize}


  \item Now for the second algebra $\F_2$, we follow a similar idea. The multiplicative identity this time is $(1,0)$ since $(1,0)\cdot (x,y)=(x,y)$, thus $\alpha(1,0)=1$. So we have a fixed value for the first basis vector, how about the second? Well consider:

    $\alpha(0,1)\cdot \alpha(0,1)=\alpha((0,1)\cdot (0,1))=\alpha(1,0)=1$

  Thus, either $\alpha(0,1)=1$ or $\alpha(0,1)=-1$, since its square is 1. For completeness let's check if both options give us a homomorphism:

  \begin{itemize}
    \item $\beta_1 (x,y)=x\beta_1(1,0)+y\beta_1(0,1)=x+y$. Evidently it is a linear map and unitality is built into it, so it remains to check whether it preserves the multiplication:
      $$\beta_1((x_1,y_1)\cdot (x_2, y_2))= \beta_1(x_1x_2+y_1y_2, x_1y_2+ x_2y_1)= x_1x_2+y_1y_2+x_1y_2+x_2y_1$$
    Which is the same as $\beta_1(x_1,y_1)\cdot\beta_1(x_2,y_2)=(x_1+y_1)(x_2+y_2)$.

    \item $\beta_2 (x,y)=x\beta_2(1,0)+y\beta_2(0,1)=x-y$. Evidently it is a linear map and unitality is built into it, so it remains to check whether it preserves the multiplication:


      $$\beta_2((x_1,y_1)\cdot (x_2, y_2))= \beta_2(x_1x_2+y_1y_2, x_1y_2+ x_2y_1)= x_1x_2+y_1y_2-x_1y_2-x_2y_1$$
    Which is the same as $\beta_2(x_1,y_1)\cdot\beta_2(x_2,y_2)=(x_1-y_1)(x_2-y_2)$.
  \end{itemize}

  Thus once again our manifold consists of the two points $M_2=\{\beta_1, \beta_2\}$, but this time a different function corresponds to an element $(x,y)\in \F_2$ namely:

    \begin{itemize}
      \item $\beta_1 \overset{(x,y)}{\mapsto} x+y$
      \item $\beta_2  \overset{(x,y)}{\mapsto} x-y$
    \end{itemize}


  \item I will finish writing up whether they are isomorphic later on.
\end{enumerate}
\end{solution}

\end{document}
